\documentclass{article}
\usepackage[utf8]{inputenc}
\usepackage{amsfonts}
\usepackage{amsmath}
\usepackage{amssymb}

\title{Correction TD Math3D}
\author{Franz Chouly}
\date{\today}

\parindent 0em

\begin{document}

\maketitle

\section{TD 1}

{\bf Exercice 1.1} --

{\bf Exercice 1.2} --

{\bf Exercice 1.3}

La fonction
$$ f(x) = 2 - x + \ln(x) $$
est définie sur $]0;+\infty[$. Sa dérivée est:
$$ f'(x) = -1 + \frac{1}{x} = \frac{1-x}{x},$$
et elle change de signe en $x=1$. Donc $f$ est croissante sur $]0;1[$ et décroissante sur $]1;+\infty[$. On a de plus:
$$ \lim_{0^+} f = -\infty, \quad f(0) = 1, \quad \lim_{+\infty} f = -\infty.$$
Comme $f$ est continue sur chacun des deux intervalles où elle est monotone, l'application du Théorème de la Bijection, adapté au cas non-borné, permet de conclure qu'il existe exactement deux zéros de $f$: $\alpha \in ]0;1[$, $\beta \in ]1;+\infty[$.
On peut avoir un encadrement de ces racines à l'aide des évaluations suivantes de la fonction $f$:
$$
f(0.1) \simeq -0.4, \quad f(0.5) \simeq 0.8, \quad f(3) \simeq 0.1, \quad f(4) \simeq -0.6, \quad f(3.5) \simeq -0.2.
$$
%>> fun(0.1)
%ans = -0.402585092994046
%>> fun(0.5)
%ans =  0.806852819440055
%>> fun(3)
%ans =  0.0986122886681098
%>> fun(4)
%ans = -0.613705638880109
%>> fun(3.5)
%ans = -0.247237031504632

\newpage

Appliquons la méthode de dichotomie pour avoir un encadrement des zéros. Voici le petit programme:
{\scriptsize
\begin{verbatim}
% Dichotomie pour le calcul de zéro

% Format long
format long;

% Précision souhaitee
precision= 1e-15;

% Valeurs initiales
x1 = 1e-6; x2 = 1;

% Fonction à considérer
fun=inline('2-x+log(x)','x');

% Calcul du zéro via Octave
zRef = fzero(fun,[x1,x2]);

disp("Zero calcule par Octave :"), disp(zRef), disp('***')

% Calcul via dichotomie
iter = 0;
while ((x2-x1)>precision)
  disp([ iter , x1 , x2 ])
  xm = 0.5*(x1+x2);
  if (fun(xm)*fun(x1) < 0)
    x2 = xm;
  else
    x1 = xm;
  end
  iter = iter + 1;
end
  
disp("*** Valeur finale :"), disp([ x1 , x2 ])
\end{verbatim}
}

%\newpage

Et les résultats, pour $\alpha$:
{\scriptsize
\begin{verbatim}
>> dichotomie
Zero calcule par Octave :
 0.158594339563039
***
   0   0.000001000000000   1.000000000000000
   1   1.00000000000e-06   5.00000500000e-01
   2   1.00000000000e-06   2.50000750000e-01
   3   0.125000875000000   0.250000750000000
   4   0.125000875000000   0.187500812500000
   5   0.156250843750000   0.187500812500000
   ...
   49  0.158594339563038   0.158594339563040
*** Valeur finale :
   0.158594339563039   0.158594339563040
\end{verbatim}
}

Et puis pour $\beta$:
{\scriptsize
\begin{verbatim}
>> dichotomie
Zero calcule par Octave :
 3.14619322062058
***
   0   3   4
   1   3.00000000000000   3.50000000000000
   2   3.00000000000000   3.25000000000000
   3   3.12500000000000   3.25000000000000
   4   3.12500000000000   3.18750000000000 <-
   ...
  49   3.14619322062058   3.14619322062058
*** Valeur finale :
   3.14619322062058   3.14619322062058    
\end{verbatim}
}

%\newpage

{\bf Exercice 1.4} 

La fonction $f$ est définie sur $\mathbb{R}$. La fonction $f$ est strictement croissante sur $\mathbb{R}$, et compte-tenu de ses limites et du fait qu'elle est continue, le Théorème de la Bijection nous permet d'assurer que $f$ n'admet qu'une seule racine $\alpha$. Comme $f(-2) = -8 -4 +1 = -11$ et $f(0)=1$, alors $\alpha$ est comprise entre $-2$ et $0$.

La méthode de Newton nous permet d'avoir une estimation plus précise:

{\scriptsize
\begin{verbatim}
>> newton
Zero calcule par Octave :
-0.453397651516404
***
   0   0
   1  -0.500000000000000
   2  -0.454545454545455
   3  -0.453398336679094
   4  -0.453397651516648
   5  -0.453397651516404
   6  -0.453397651516404
   ...
*** Valeur finale de xR :
-0.453397651516404
\end{verbatim}
}

Et voici le script associé:
{\scriptsize
\begin{verbatim}
% Newton pour le calcul de zéro

% Format long
format long;

% Nombre d'itérations de Newton
N=20;

% Valeur initiale
xI = 0;

% Fonction à considérer
fun=inline('x^3+2*x+1','x');
% Derivee de la fonction
fprime=inline('3*x^2+2','x');

% Calcul du zéro via Octave
zRef = fzero(fun,[-2,0]);

disp("Zero calcule par Octave :"), disp(zRef), disp('***')

% Calcul via Newton
xR = xI;
for i=1:N 
  disp([ i-1 , xR ])
  xR = xR - fun(xR)/fprime(xR);
end
  
disp("*** Valeur finale de xR :"), disp(xR)
\end{verbatim}
}

{\bf Exercice 1.5}

Il suffit de considérer $f(x) = x - g(x)$, et d'exploiter l'inclusion $g([a,b]) \subset [a,b]$, qui donne $f(a) \leq 0$ et $f(b) \geq 0$.

\newpage

\section{TD 2}

{\bf Exercice 2.1}

\begin{enumerate}
    \item On calcule:
    \[ f(1) = e - 2 - 1 = e - 3 < 0, \quad f(2) = e^2 - 2*2 - 1 = e^2 - 5 > 0 \]
    et comme f est continue sur $\mathbb{R}$, le TVI nous assure l'existence (d'au moins) un zéro $\xi$ dans l'intervalle $[1,2]$.
    
    \item Il suffit de partir de l'expression
    $$  e^x - 2x - 1 = 0$$ qui devient $$ e^x = 2 x + 1,$$ et comme $x\in[1,2]$, on peut prendre le logarithme de cette dernière expression.
    
    \item La fonction $g$ est continue sur $[1,2]$, elle est croissante et de plus $g(1) = \ln(3) \geq 1$, $g(2) = \ln(5)\leq2$. On peut donc appliquer le Théorème de Brouwer et déduire que $g$ admet au moins un point fixe sur $[1,2]$.
    
    \item Non: $g$ ne vérifie plus la condition $g([1,2]) \subset [1,2]$.
\end{enumerate}

{\bf Exercice 2.2}

\begin{enumerate}
    \item $g$ est Lipschitz, donc continue.
    \item Le Théorème de Brouwer nous garantit qu'on va avoir au moins un point fixe. Il est forcément unique. En effet, soient $\xi_1$ et $\xi_2$ deux points fixes:
    \[
    | \xi_1 - \xi_2 | = | g(\xi_1) - g(\xi_2) |
    \leq L | \xi_1 - \xi_2 |
    \]
    où on a utilisé la définition d'un point fixe et la propriété contractante de $g$. La relation ci-dessus donne alors:
    \[
    (1 - L) | \xi_1 - \xi_2 | \leq 0
    \]
    qui implique $\xi_1 = \xi_2$.
    \item On a 
    \[
    | x_{k+1} - \xi |
    = |g(x_k) - g(\xi) | \leq L | x_k - \xi |
    \]
    en utlisant d'abord la relation de recurrence et le fait que $\xi$ est un point fixe, puis la propriété contractante de $g$.
    Une récurrence immédiate donne alors:
    \[
    | x_k - \xi | \leq L^k | x_0 - \xi |.
    \]
    Comme $0 < L < 1$, ceci implique que $(x_k)$ converge vers $\xi$.
\end{enumerate}

{\bf Exercice 2.3}

\begin{enumerate}
    \item Il s'agit d'une application du Théorème des Accroissements Finis.
    
    \item On a $g'(x) = \frac{2}{2x+1}$ qui est strictement décroissante (composition d'une fonction croissante par une fonction décroissante), sur $[1,2]$.
    
    On a donc $$ g'(2) \leq g'(\eta) \leq g'(1) $$
    et $g'(2) = \frac{2}{5}$, $g'(1) = \frac{2}{3}$.
    
    \item C'est une conséquence directe des deux questions précédentes, et on a $L = 2/3$.
    
    \item Il suffit d'appliquer le Théorème du Point Fixe de Banach.
    
     \item On obtient les valeurs suivantes :

\begin{verbatim}
0  1.000000
1  1.098612
2  1.162283
3  1.201339
4  1.224563
5  1.238121
6  1.245952
7  1.250447
8  1.253018
9  1.254486
10 1.255323
11 1.255800
\end{verbatim}
    
\end{enumerate}

{\bf Exercice 2.4}

\begin{enumerate}
    \item On peut par exemple partir de l'inégalité triangulaire 
    \[
    | x_0 - \xi | \leq | x_0 - x_1 | + | \xi - x_1 |
    \]
    et puis ensuite utiliser
    \[
    | \xi - x_1 | = | g(\xi) - g(x_0) | \leq L 
    | \xi - x_0 |
    \]
    pour obtenir
    \[
    | x_0 - \xi | \leq | x_0 - x_1 | +
    L | \xi - x_0 |,
    \]
    ce qui donne bien l'inégalité voulue.
    \item On obtient le résultat en combinant la question 1 et la propriété de contraction de $g$, puis en faisant une récurrence immédiate.
    
    \item On part de la relation obtenue à la question 2 et on écrit:
    \begin{eqnarray*}
    (1-L)^{-1} L^k |x_1 - x_0| & \leq & \varepsilon\\
    \ln((1-L)^{-1} ) + k \ln(L) + \ln |x_1 - x_0| & \leq & \ln \varepsilon \\
    k \ln(L) + \ln |x_1 - x_0| & \leq & \ln \varepsilon - \ln((1-L)^{-1} ) \\
    k \ln(L) + \ln |x_1 - x_0| & \leq & \ln ((1-L)\varepsilon)
    \end{eqnarray*}
    d'où le résultat.
    
    \item Il faut 33 itérations.
    
\end{enumerate}

\newpage

\section{TD 3}

{\bf Exercice 3.1}

Nous allons montrer que $g$ est contractante dans un voisinage du point fixe $\xi$. Comme la dérivée de $g$ est continue dans un voisinage de $\xi$, elle l'est en particulier en $\xi$, et cela implique que pour tout $\varepsilon > 0$, il existe $\delta > 0$ tel que pour tout $x\in \mathbb{R}$ qui vérifie $|x - \xi| < \delta$, nous avons $x\in I$ et:
\[
| g'(x) - g'(\xi) | < \varepsilon. 
\]
Définissons maintenant
\[
0 \leq L = \frac12 ( 1 + |g'(\xi)| ) < 1
\]
et
\[
\varepsilon = \frac12 ( 1 - |g'(\xi)| ) > 0.
\]
Alors pour tout $x \in I$ tel que $|x - \xi| < \delta$, nous avons:
\[
|g'(x)| \leq |g'(\xi)| + |g'(x) - g'(\xi)| 
\leq |g'(\xi)| + \varepsilon = L.
\]
On a $]\xi - \delta ; \xi + \delta [ \subset I$.
En appliquant le Théorème des Acroissements Finis, on voit alors que $g$ est contractante sur $]\xi - \delta ; \xi + \delta [$. Il suffit alors de reprendre le même raisonnement que dans la démonstration du Théorème de Point Fixe de Banach-Picard pour montrer que la suite ($x_k$) donnée par la relation de récurrence $x_{k+1} = g(x_k)$ converge vers $\xi$, si $x_0$ est dans l'intervalle $I \cap ]\xi - \delta ; \xi + \delta [$ (on vérifie en particulier qu'à chaque itération on reste bien dans le bon intervalle $]\xi - \delta ; \xi + \delta [$).

\medskip

{\bf Exercice 3.2}

Il suffit tout simplement de reconnaître le taux d'accroissement
\[
\frac{|g(x_k)-g(\xi)|}{|x_k - \xi|}
\]
qui tend vers $g'(\xi)$ lorsque $k$ tend vers l'infini.

\medskip

{\bf Exercice 3.3}

\begin{enumerate}
    \item Lorsque $x \rightarrow 0+$:
    \begin{eqnarray*}
    \log_2 (1/x) & \rightarrow & +\infty\\
    (\log_2 (1/x))^{\frac{1}{\alpha}} & \rightarrow & +\infty
    \end{eqnarray*}

    Et on voit alors que la limite de $g$ en $0+$ est $0$. On a alors $g$ qui est continue sur $[0,1]$.

\item On vérifie directement que $g$ est croissante sur $[0,1]$.

\item Le seul point fixe de $g$ est $0$. En effet, on a
\[
g(0) = 0, \quad g(x) < x.
\]
On obtient la deuxième relation en partant de l'inégalité $\log_2 (1/x) < 1 + \log_2 (1/x)$.

\item Il suffit de supposer qu'au rang $k$, on a $x_k = 2^{-k^\alpha} > 0$ et on calcule $x_{k+1}$:
\begin{eqnarray*}
x_{k+1} & = & g(x_k)\\
& = & 2^{-(1+\log_2(1/x_k)^{\frac{1}{\alpha}})^\alpha}\\
& = & 2^{-(1+\log_2(2^{k^\alpha})^{\frac{1}{\alpha}})^\alpha}\\
& = & 2^{-(1+k)^\alpha},
\end{eqnarray*}
ce qui permet de conclure: la suite $(x_k)$ converge donc vers $0$, unique point fixe de $g$.

\item On calcule:
\[
\frac{x_{k+1}}{x_k} = 2^{-(1+k)^\alpha + k^\alpha }
 = 2^{\left (1 -\left (1+\frac{1}{k} \right )^\alpha \right ) k^\alpha}
\]
et on utilise l'équivalent
\[
1 -\left (1+\frac{1}{k} \right)^\alpha \sim -\frac{\alpha}{k} 
\]
pour conclure
\[
\frac{x_{k+1}}{x_k} \sim 2^{-\alpha k^{\alpha-1}}
\]
et en déduire que si $\alpha < 1$, alors la limite est 1 (convergence sous-linéaire), si $\alpha = 1$, la limite est $\frac12$ (convergence linéaire), et si $\alpha > 1$, la limite est 0 (convergence super-linéaire).

{\bf Bonus:} On peut regarder concrètement les différences de comportement sous-linéaire/linéaire/super-linéaire en calculant explicitement les premiers termes de la suite $x_k$. Le programme Octave qui permet cela est donné ci-dessous:

\begin{verbatim}
% Suite x_k de l'exercice 3.3

alpha1 = 2;
alpha2 = 1;
alpha3 = 0.5;

k = 1:10;

xk = [ 2.^(-k.^alpha1)' 2.^(-k.^alpha2)' 2.^(-k.^alpha3)' ]
\end{verbatim}

On obtient alors pour les premières valeurs $x_0,\ldots,x_9$:

\begin{verbatim}
   alpha = 2    alpha = 1    alpha = 0.5

   5.0000e-01   5.0000e-01   5.0000e-01
   6.2500e-02   2.5000e-01   3.7521e-01
   1.9531e-03   1.2500e-01   3.0102e-01
   1.5259e-05   6.2500e-02   2.5000e-01
   2.9802e-08   3.1250e-02   2.1226e-01
   1.4552e-11   1.5625e-02   1.8308e-01
   1.7764e-15   7.8125e-03   1.5979e-01
   5.4210e-20   3.9062e-03   1.4079e-01
   4.1359e-25   1.9531e-03   1.2500e-01
   7.8886e-31   9.7656e-04   1.1170e-01
\end{verbatim}

Le comportement super-linéaire se manifeste bien pour $\alpha=2$ où la convergence de $(x_k)$ vers le point fixe $\xi = 0$ est extr\^emement rapide. Pour $\alpha=1$, on retrouve exactement le comportement linéaire, avec le facteur $\frac12$ qui donne le ratio entre deux itérations successives. Pour $\alpha=0.5$, le comportement se manifeste effectivement par une convergence très lente.

\end{enumerate}

{\bf Exercice 3.4 (correction esquissée)}

Déjà, on montre que $g$ est ``dilatante" dans un voisinage $I_\delta$ de $\xi$, de la même façon qu'à l'exercice 3.2: on a donc l'existence de $L>1$ tel que pour tout $x$ et tout $y$ dans $I_\delta$:
\[
|g(x)-g(y)| \geq L |x-y| \qquad (*)
\]
Ensuite considérons la suite $(x_k)$, et plaçons nous dans le cas de figure où $\xi$ n'est pas atteint après un nombre fini d'itérations. Supposons qu'à une itération donnée $k_0$, $x_{k_0}$ appartienne à $I_\delta$. Alors compte-tenu de $(*)$ il y aura un rang $k > k_0$ tel que $|x_k - \xi | > \delta$.

Donc il existe $\delta > 0$ tel que pour tout $k_0$, il existe $k\geq k_0$ qui vérifie $x_k \notin I_\delta$: la suite ($x_k$) ne converge pas vers $\xi$.

\medskip

{\bf Exercice 3.4 (correction détaillée)}

\smallskip

{\bf Etape 1:} montrons que $g$ est `dilatante'.

\smallskip

Comme pour l'exercice 3.1, on choisit

\[
\varepsilon = \frac12 (|g'(\xi)| - 1), \qquad L = \frac12 ( |g'(\xi)| + 1 ).
\]
L'hypothèse $|g'(\xi)|>1$ nous assure alors que $\varepsilon > 0$ et $L>1$.

La continuité de $g'$ en $\xi$ nous permet alors d'affirmer qu'il existe $\delta > 0$ tel que pour tout $x$ dans $I_\delta = [ \xi - \delta ; \xi + \delta ] \subset I$, nous avons:
\[
|g'(x) - g'(\xi)| \leq \varepsilon.
\]
Alors pour $x \in I_\delta$, nous appliquons l'inégalité triangulaire:
\[
|g'(\xi)| \leq |g'(\xi) - g'(x)| + |g'(x)|,
\]
ce qui donne la minoration:
\[
|g'(x)| \geq |g'(\xi)| - |g'(\xi) - g'(x)| \geq |g'(\xi)| - \varepsilon = L.
\]
Résumons:
\[
|g'(x)| \geq  L > 1, \qquad \forall x \in I_\delta.
\]
Le Théorème des accroissements finis nous permet alors de conclure que pour tout $x$ et tout $y$ dans $I_\delta$:
\[
|g(x)-g(y)| \geq L |x-y| \qquad (*)
\]

\smallskip

{\bf Etape 2:} on va montrer que toute suite de point fixe qui n'atteint pas $\xi$ en un nombre fini d'itérations ne converge jamais vers $\xi$.

\smallskip

Supposons que la suite $(x_k)$ donnée dans l'énoncé n'atteint jamais $\xi$ en un nombre fini d'itérations.

\smallskip

On peut déjà traiter un premier cas où chaque terme $x_k$ n'appartient pas à $I_\delta$, pour toutes les valeurs de $k$: dans ce cas on a bien vérifié que la suite $(x_k)$ ne converge pas vers $\xi$. 

\smallskip

Maintenant, on va supposer qu'à une itération donnée $k_0 \geq 0$, $x_{k_0}$ appartienne à $I_\delta$. Alors compte-tenu de $(*)$ il y aura un rang $k > k_0$ tel que $|x_k - \xi | > \delta$.


On peut s'en assurer avec un raisonnement par l'absurde: si ce n'était pas le cas, on aurait, pour tout $k \geq k_0$, $|x_k - \xi| \leq \delta$, ce qui signifie que $x_k \in I_\delta$. On peut alors appliquer $(*)$, et une récurrence immédiate donne, pour tout $k \geq k_0$:
\[
|x_{k} -\xi | \geq L^{k-k_0} |x_{k_0} - \xi|.
\]
Comme on est dans le cas où $x_{k_0} \neq \xi$, et comme $L>1$, cela signifie que $|x_{k} -\xi |$ tend vers l'infini, ce qui est contradictoire avec l'hypothèse $x_k \in I_\delta$.

\medskip

Récapitulons et reformulons alors ce qu'on vient de montrer : 

\smallskip

Il existe $\delta > 0$ tel que pour tout $k_0$, il existe $k\geq k_0$ qui vérifie $|x_k - \xi | > \delta$.

\smallskip

C'est exactement la négation de la définition de la convergence d'une suite vers la limite finie $\xi$. On a donc bien que la la suite ($x_k$) ne converge pas vers $\xi$.

\medskip

{\bf Exercice 3.5}

\begin{enumerate}
    \item En résolvant $x^2 - 2x + c = 0$, on trouve tous les points fixes de $g$.
    On a
    \[
    \Delta = 4 ( 1 - c )
    \]
    et si $c<1$, on a deux points fixes
    \[
    \xi_1 = 1 - \sqrt{1-c}, \quad
    \xi_2 = 1 + \sqrt{1-c}.
    \]
    Si $c=1$, $1$ est le seul point fixe, et si $c>1$, il n'y a pas de point fixe.
    
    \item --
    
    \item D'abord calculons $g'(x)=x$. On a
    \[
    |g'(\xi_2)| = |\xi_2| > 1
    \]
    ce qui montre que $\xi_2$ est instable.
    
    \item On a d'un côté $g'(\xi_1) < 1$ pour $c<1$ et
    \[
    g'(\xi_1) = 1 - \sqrt{1-c} > -1
    \]
    pour $\sqrt{1-c}<2$ et donc $(1-c)<4$: il faut donc que $c>-3$ pour que $\xi_1$ soit stable.
    
    \item On a
    \[
    x_{k+1} - \xi_1 = g(x_k) - g(\xi_1) = \frac12(x_k^2 + c ) - \frac12(\xi_1^2 + c) = 
    \frac12(x_k^2 - \xi^2_1 )
    \]
    d'où
    \[
    x_{k+1} - \xi_1 = \frac12 (x_k + \xi_1 )(x_k-\xi_1).
    \]
    Soit $\delta > 0$ tel que $0\leq x_0 < \xi_2 - \delta$.
    Supposons que $0\leq x_k < \xi_2 - \delta$, alors
    \[
    0 \leq x_k + \xi_1 < \xi_2 + \xi_1 - \delta < 2 - \delta.
    \]
    Donc
    \[
    |x_{k+1} - \xi_1| \leq \frac12 |x_k + \xi_1 ||x_k-\xi_1| \leq \frac12(2-\delta) |x_k-\xi_1|.
    \]
    On a donc via une récurrence immédiate:
    \[
    |x_{k} - \xi_1| \leq \left (1-\frac\delta2 \right)^k |x_0-\xi_1|
    \]
    ce qui permet de conclure.
    \item
    La vitesse de convergence est donnée par
    \[
    g'(\xi_1) = 1 - \sqrt{1-c}
    \]
    et donc lorsque $c=0$ la convergence est sur-linéaire, et lorsque $c$ est proche de 1 la convergence est sous-linéaire.
\end{enumerate}

\newpage

\section{TD 4}

{\bf Exercice 4.1}

Soit $g$ la fonction de point fixe associée à la méthode de relaxation.
Compte-tenu des hypothèses sur $f$, $g$ est définie et continue dans un voisinage de $\xi$, et aussi dérivable et de dérivée continue dans un voisinage de $\xi$.
Comme $g(x) = x - \lambda f(x)$, nous avons:
\[
g'(x) = 1 - \lambda f'(x).
\]
Pour pouvoir appliquer le Théorème 3 (Convergence pour un point fixe stable) vu au TD 3, ou Théorème d'Ostrowski, il faut,
en $x=\xi$, vérifier $|g'(\xi)| < 1$, ce qui est équivalent à dire:
\[
-1 < g'(\xi) < 1.
\]
Reformulons cette condition:
$$
-1< 1  - \lambda f'(x) < 1
$$
et donc on obtient
$$
-2< - \lambda f'(x) < 0
$$
autrement dit
\[
0 < \lambda f'(\xi) < 2,
\]
ce qui est toujours possible en prenant $\lambda$ du même signe que $f'(\xi)$ et
\[
|\lambda| < 2/|f'(\xi)|.
\]
Donc pour $\lambda$ 'suffisamment petit' en valeur absolue (inférieur à la constante ci-dessus), et du bon signe, l'application du Théorème 3 nous garantit que la suite $(x_k)$ donnée par la méthode de relaxation (1) va converger vers le zéro $\xi$, pourvu que $x_0$ soit suffisamment proche de $\xi$.

{\bf Remarque :} La vitesse de convergence ici est donnée par :
\[
|g'(\xi)| = | 1 - \lambda f'(\xi) |
\]
En particulier, on aura une convergence super-linéaire pour 
\[
\lambda = \frac{1}{f'(\xi)}
\]
et on aura une convergence sous-linéaire pour
\[
\lambda \simeq 0, \qquad \lambda \simeq \frac{2}{f'(\xi)}.
\]
Le résultat démontré dans l'exercice 3.4 permet aussi de voir que pour
\[
|\lambda| > 2/|f'(\xi)|,
\]
la suite de relaxation ne peut plus converger vers $\xi$.

\medskip

{\bf Exercice 4.2}

\begin{enumerate}
    \item En passant à la limite quand $k \rightarrow +\infty$ dans l'équation 
    $$
    x_{k+1} = x_k - \lambda(x_k)f(x_k),
    $$
    on aboutit a 
    \[
    \xi = \xi - \lambda(\xi) f(\xi).
    \]
    Comme $\lambda(\xi)\neq0$ alors nécessairement $f(\xi)=0$.
    
    \item 
    On va supposer que $\lambda$ est dérivable dans un voisinage de $\xi$.
    La vitesse de convergence asymptotique est donnée par $|g'(\xi)|$ (voir Exercice 3.2) et
    \[
    g'(\xi) = 1 - \lambda'(\xi) f(\xi) - \lambda(\xi) f'(\xi) = 1 -\lambda(\xi)f'(\xi),
    \]
    car $f(\xi)=0$.
    
    \item On voit que pour avoir une convergence super-linéaire, on peut prendre $\lambda(x) = 1/f'(x)$.
\end{enumerate}

{\bf Exercice 4.3}

\begin{enumerate}
    \item On écrit la formule de Taylor-Lagrange au point $x_k$ et à l'ordre 1: il existe $\eta_k \in I_\delta$ ($|\eta_k - \xi| \leq |x_k - \xi| \leq h \leq \delta$) qui vérifie
    \[
    f(\xi) = f(x_k + (\xi - x_k)) = f(x_k) + (\xi - x_k)f'(x_k) + \frac12 (\xi - x_k)^2 f''(\eta_k).
    \]
    On conclut avec l'égalité $f(\xi)=0$:
    \[
    0 = f(\xi)  = f(x_k) + (\xi - x_k)f'(x_k) + \frac12 (\xi - x_k)^2 f''(\eta_k).
    \]
    
    
    \item On utilise la définition de la suite donnée par la méthode de Newton:
    \begin{eqnarray*}
    \xi - x_{k+1} &=& \xi - (x_k - f(x_k)/f'(x_k) )\\
    &=& \frac{(\xi-x_k)f'(x_k) + f(x_k)}{f'(x_k)},
    \end{eqnarray*}
    puis on reconnait au numérateur une partie de l'expression donnée à la question 1. 
    On obtient bien alors:
    \[
    \xi - x_{k+1} = - \frac{(\xi - x_k)^2 f''(\eta_k)}{2 f'(x_k)}.
    \]
    Puis on majore cette dernière expression:
    \[
    | \xi - x_{k+1} | \leq \frac12 |\xi - x_k| \frac{|f''(\eta_k)|}{|f'(x_k)|} |\xi - x_k|.
    \]
    Puis on utilise $|\xi - x_k| \leq h$ et $\frac{|f''(\eta_k)|}{|f'(x_k)|} \leq A \leq 1/h$ pour conclure:
    \[
    | \xi - x_{k+1} | \leq \frac12 h \frac{1}{h} |\xi - x_k| \leq \frac12 |\xi - x_k|.
    \]
    \item On obtient le résultat voulu par récurrence immédiate:
    
    Au rang 1:
    \[
    | \xi - x_{1} |  \leq \frac12 |\xi - x_0| \leq \frac12 h.
    \]

    Au rang $k+1$:
    \[
    | \xi - x_{k+1} |  \leq \frac12 |\xi - x_k| \leq \frac12 2^{-k} h
    \leq 2^{-(k+1)} h.
    \]
    
    Et on en déduit qu'il y a convergence de $(x_k)$ vers $\xi$.
    
    \item Il faut partir de :
     \[
    \xi - x_{k+1} = - \frac{(\xi - x_k)^2 f''(\eta_k)}{2 f'(x_k)}.
    \]
    On peut transformer cette expression comme suit:
     \[
    \frac{\xi - x_{k+1}}{(\xi - x_k)^2} = - \frac{f''(\eta_k)}{2 f'(x_k)}.
    \]
    On prend ensuite les valeurs absolues:
     \[
    \left | \frac{\xi - x_{k+1}}{(\xi - x_k)^2} \right | = \frac12 \left |\frac{f''(\eta_k)}{f'(x_k)} \right |.
    \]
    A la limite quand $k$ tend vers l'infini, on a $x_k$ qui tend vers $\xi$, et $\eta_k$ aussi (car $|\eta_k-\xi|\leq|x_k -\xi|$). Comme $f'$ et $f''$ sont continues,
    en passant à la limite, on obtient finalement $\mu = \frac12|f''(\xi)/f'(\xi)|$ qui est bien plus petit que $A/2$.
    
    \medskip

    On vient ici de montrer le Théorème de Kantorovich, qui établi la convergence quadratique locale de la méthode de Newton.    
\end{enumerate}

{\bf Exercice 4.4}

\medskip

Supposons que $x_k \in J = [ \xi , X ]$. On a la relation suivante pour les itérations de Newton:
\[
x_{k+1} = x_k - \frac{f(x_k)}{f'(x_k)}. 
\]
Cette relation est bien définie car $f'(x_k) > 0$, d'après l'énoncé.
L'énoncé nous affirme de plus que $f'>0$ sur $J$, comme $f(\xi)=0$, cela signifie aussi que $f\geq0$ sur $J$. On déduit alors de la relation ci-dessus:
\[
x_{k+1} \leq x_k \qquad (P_1)
\]
Souvenons nous alors de la relation trouvée à l'exercice précédent :
\[
    \xi - x_{k+1} = - \frac{(\xi - x_k)^2 f''(\eta_k)}{2 f'(x_k)}.
\]
On a $\eta_k \in [\xi;x_k] \subset J$, donc $f''(\eta_k)> 0$. De plus $f'(x_k)>0$. Ces deux propriétés combinées permettent d'affirmer que
\[
    - \frac{(\xi - x_k)^2 f''(\eta_k)}{2 f'(x_k)} \leq 0 
\]
et donc
\[
    \xi - x_{k+1} \leq 0,
\]
autrement dit :
\[
\xi \leq  x_{k+1} \qquad (P_2)
\]
Reste à combiner $(P_1)$ et $(P_2)$. Ces deux propriétés, combinées à une récurrence immédiate (valable car $x_0 \in J$) permettent d'avoir, pour tout $k\geq0$:
\[
\xi \leq  x_{k+1} \leq x_k.
\]
On vient donc de montrer que la suite $(x_k)$ est décroissante et minorée. Elle converge donc vers une limite $x^*$.

Comme $f$ et $f'$ sont continues sur $J$, on peut passer à la limite dans la relation de récurrence:
\[
x_{k+1} = x_k - \frac{f(x_k)}{f'(x_k)}. 
\]
Cela donne:
\[
x^* = x^* - \frac{f(x^*)}{f'(x^*)}. 
\]
Comme $f'(x^*)>0$, ceci implique $f(x^*)=0$, or, comme $f(x)>0$ si $x>\xi$, on a alors nécessairement $x^*=\xi$.

On obtient un taux quadratique de convergence exactement comme à l'exercice précédent.

\medskip

On a donc bien montré que pour tout $x_0 \in J$, la suite $(x_k)$ donnée par les itérées de Newton converge vers $\xi$, unique zéro de $f$ sur $J$.

\medskip

{\bf Exercice 4.5}

\begin{enumerate}
    \item Donnons déjà le programme Octave associé à la méthode de relaxation:

{\footnotesize
\begin{verbatim}
% Relaxation pour le calcul de zéro

% Format long
format long;

% Nombre d'itérations de relaxation
N=20;

% Paramètre de relaxation
lambda=1e-1;

% Valeur initiale
xI = 0;

% Fonction à considérer
fun=inline('x^3+2*x+1','x');

% Calcul du zéro via Octave
zRef = fzero(fun,[-2,0]);

disp("Zero calcule par Octave :"), disp(zRef), disp('***')

% Calcul via relaxation
xR = xI;
for i=1:N 
  disp([ i-1 , xR ])
  xR = xR - lambda * fun(xR);
end
  
disp("*** Valeur finale de xR :"), disp(xR)
\end{verbatim}
}

Et faisons déjà un test pour $\lambda = 0.1$ (on donne sur la colonne de gauche l'itération $i$ et sur la colonne de droite la valeur de $x_i$ correspondante):

{\footnotesize
\begin{verbatim}
   0    0
   1   -0.100000000000000
   2   -0.179900000000000
   3   -0.243337771460100
   ...
   18  -0.450846888893672
   19  -0.451513465729099
   20  -0.452006019965466
\end{verbatim}
}

La valeur précise du zéro (calculée avec la fonction dédiée d'Octave) est %Zero calcule par Octave :
$$-0.453397651516404$$ On obtient donc au bout de 20 itérations une précision de 2 chiffres après la virgule.

En augmentant la valeur de $\lambda$ jusqu'à $0.7$ environ, on obtient un comportement identique, mais avec une convergence plus rapide vers la solution. Par contre, pour $\lambda=0.8$, on obtient le comportement suivant :

{\footnotesize
\begin{verbatim}
   0   0
   1  -0.800000000000000
   2   0.0896000000000001
   3  -0.854335458508800
   4   0.211457369600640
   5  -0.934438542693973
   6   0.413406116828508
   7  -1.10456608148935
   8   0.940854667322071
   9  -2.03079409051181
   10  7.11867480910851
   11 -293.6653058752566
   12  2.02605705193333e+07
   ...
\end{verbatim}
}

On se rend compte que la méthode ne converge plus du tout.

D'un autre côté, si on prend $\lambda = 0.001$, on obtient:

{\footnotesize
\begin{verbatim}
   0   0
   1  -0.00100000000000000
   2  -0.00199799999900000
   ...
   19 -0.0186618172987305
   20 -0.0196244871649046
\end{verbatim}
}

On se rapproche (symboliquement ...) de la valeur attendue, mais avec une lenteur extrême : en 20 itérations, nous n'avons même pas une approximation à la première décimale.

{\bf En conclusion :} on retrouve bien ici le comportement théorique attendu et prédit à l'Exercice 4.1, autrement dit que la méthode converge à condition que $\lambda$ soit suffisamment petit. On voit de plus que la méthode de relaxation est susceptible de diverger pour des valeurs trop grandes de $\lambda$, autrement dit la condition "$\lambda$ suffisamment petit", en plus d'être suffisante, semble aussi nécessaire. On voit également que, d'un point de vue pratique, choisir $\lambda$ trop petit permet, certes, d'assurer la convergence de la méthode, mais ne permet pas d'aboutir à une approximation acceptable même au bout de 20 itérations.

Bref, une des difficultés pratiques associée à la méthode de relaxation est qu'il est toujours délicat de bien choisir $\lambda$.

\medskip

{\bf Remarque :} un calcul approché de la dérivée de $f$ en $\alpha$ donne
\[
f'(\alpha) \simeq  2.61670829120177
\]
ce qui permet d'obtenir une borne théorique pour le choix de $\lambda$. En effet, on doit avoir, d'après l'Exercice 4.1:
\[
0 \leq \lambda \leq \frac{2}{f'(\alpha)}
\]
avec $2/f'(\alpha) \simeq 0.77$. Noter que d'après les résultats (Exercice 3.1 et Exercice 3.4) du TD3, on doit avoir convergence de la méthode pour $\lambda$ inférieur à ce seuil (pour $x_0$ suffisamment proche de $\alpha$), \emph{et} divergence de la méthode pour $\lambda$ supérieur à ce seuil. C'est effectivement ce qu'on observe ici en pratique.

{\bf Remarque (bis) :} remarquer qu'en plus, on aurait pu étudier l'influence du choix de la valeur initiale $x_0$.

\item Donnons déjà le programme Octave associé à la méthode de Newton :

{\footnotesize 
\begin{verbatim}
% Newton pour le calcul de zéro

% Format long
format long;

% Nombre d'itérations de Newton
N=5;

% Valeur initiale
xI = 0;

% Fonction à considérer
fun=inline('x^3+2*x+1','x');
% Derivee de la fonction
fprime=inline('3*x^2+2','x');

% Calcul du zéro via Octave
zRef = fzero(fun,[-2,0]);

disp("Zero calcule par Octave :"), disp(zRef), disp('***')

% Calcul via Newton
xR  = xI;
xRt = [ xR ];
cvFact = 0;
for i=1:N 
  disp([ i-1 , xR, cvFact ])
  xR = xR - fun(xR)/fprime(xR);
  xRt = [ xRt xR ];
  cvFact = abs(xRt(i+1) - zRef)/abs((xRt(i) - zRef)^2);
end
  
disp("*** Valeur finale de xR :"), disp(xR)
\end{verbatim}
}

On obtient le résultat suivant (première colonne : indice de l'itération $i$ ; deuxième colonne : valeur de $x_i$ ; troisième colonne : taux de convergence quadratique empirique ):

\begin{verbatim}
   0   0                   0
   1  -0.500000000000000   0.226698825758202
   2  -0.454545454545455   0.528508236915050
   3  -0.453398336679094   0.520066611476640
   4  -0.453397651516648   0.519817033317573
*** Valeur finale de xR :
-0.453397651516404
***
\end{verbatim}

On aboutit en seulement 5 itérations à un résultat aussi précis que la machine est capable de les délivrer : 15 chiffres significatifs exacts après la virgule. C'est incomparablement mieux que la méthode de relaxation.

\item On vérifie bien, via le calcul du taux de convergence empirique, qu'on a une convergence quadratique, avec un taux de convergence associé d'environ 0.52.

En effet, si on appelle 
\[
e_k = | x_k - \xi |
\]
l'erreur d'approximation au rang $k$. On a vu à l'exercice 4.3 que la quantité
\[
\frac{e_{k+1}}{e_{k}^2}
\]
tendait vers une limite finie (c'est ça, la vitesse de convergence quadratique).

\end{enumerate}

{\bf Exercice 4.6}

On trouve approximativement: $r\simeq 0.06140241153618$.

\newpage

\section{TD 5}

{\bf Exercice 5.1}

\begin{enumerate}
    \item Il suffit de construire explicitement les polynômes souhaités ($L_0,\ldots,L_n$). Soit $k$ un entier compris entre $0$ et $n$. On remarque déjà que le polynôme
    \[
    (x - x_0) ( x - x_1) \ldots (x - x_{k-1}) (x - x_{k+1}) \ldots (x - x_n)
    \]
    est un polynôme de degré $n$ qui s'annule pour tous les points $x_i$ tels que $i$ est différent de $k$. Pour vérifier la condition $L_k(x_k)=1$, il suffit alors de prendre:
    \[
    L_k(x) = \frac{(x - x_0) ( x - x_1) \ldots (x - x_{k-1}) (x - x_{k+1}) \ldots (x - x_n)}{(x_k - x_0) ( x_k - x_1) \ldots (x_k - x_{k-1}) (x_k - x_{k+1}) \ldots (x_k - x_n)}.
    \]
    On vérifie bien que le dénominateur ne s'annule pas, car les réels $x_i$ sont tous distincts.
    
    \item Le polynôme 
    \[
    p_n(x) = \sum_{k=0}^n y_k L_k(x)
    \]
    est bien un polynôme de degré maximal $n$. De plus, il 
    vérifie bien la condition (4). En effet
    \[
    p_n(x_i) = \sum_{k=0}^n y_k L_k(x_i) = y_i L_i(x_i) = y_i.
    \]
    Donc : 
    \[
    p_n(x_i)=y_i, \qquad i=0,\ldots,n.
    \]
    
    \item Prenons deux polyn\^omes $q_n$ et $p_n$ qui vérifient la condition (4). Alors le polynôme $p_n-q_n$ admet exactement $n+1$ racines distinctes. Comme c'est un polynôme de degré $n$, il s'agit alors nécessairement du polynôme nul. En conséquence :
    \[
    p_n = q_n.
    \]
    
    \item Lorsque $n=0$, le problème d'interpolation de Lagrange se réduit à savoir s'il existe un unique polynôme $p_0$ de degré maximal $0$ dont le graphe associé passe par un point de coordonnées $(x_0,y_0)$. Il suffit alors de prendre: $p_0 = y_0$.
\end{enumerate}

{\bf Remarque :} si on regarde les choses du point de vue de l'algèbre linéaire, on a montré alors que $(L_k)$ est une base de l'espace vectoriel $\mathcal{P}_n$.

\medskip

Soit maintenant, $\mathcal{L}(\mathcal{P}_n;\mathbb{R})$, l'espace vectoriel dual de $\mathcal{P}_n$ (espace vectoriel des formes linéaires sur $\mathcal{P}_n$), et soit
$$
\psi_j : p \rightarrow p(x_j),
$$
alors $(\psi_j)$ est une base de $\mathcal{L}(\mathcal{P}_n;\mathbb{R})$, et c'est la base duale de $(L_k)$.

\medskip

Introduisons finalement
\[
\Psi : \mathcal{P}_n \rightarrow \mathbb{R}^{n+1}
\]
définie par
\[
\Psi(p) = ( p(x_0), \ldots , p(x_n) ).
\]
On a montré que $\Psi$ est injective, donc bijective.

\bigskip

{\bf Exercice 5.2}

On explicite d'abord dans ce cas particulier les polynômes de la base de Lagrange, ici de degré 2 :
\[
L_0(x) = \frac12 x ( x - 1 ), \qquad
L_1(x) = 1 - x^2, \qquad
L_2(x) = \frac12 x ( x + 1 ).
\]
Il suffit ensuite d'écrire :
\[
p_2 (x) = e^{-1} \frac12 x(x-1) + e^0(1-x^2) + e^1 \frac12 x(x+1).
\]
En ré-ordonnant les termes et en simplifiant, on obtient finalement :
\[
p_2 (x) = 1 + (\sinh{1}) x + (\cosh{1} - 1) x^2.
\]

{\bf Exercice 5.3}

Soit 
\[
\varphi(t) = f(t) - p_n(t) - \frac{f(x)-p_n(x)}{\pi_{n+1}(x)} \pi_{n+1} (t).
\]

\begin{enumerate}
    \item La fonction $\varphi$ est bien définie sur l'intervalle $[a,b]$ car $\pi_{n+1}(x)$ ne s'annule jamais ($x$ est distinct de tous les $x_i$) et c'est alors une combinaison linéaire d'une fonction $f$, bien définie par hypothèse, et de polynômes.
    
    \item La fonction $\varphi$ s'annule en tous les points $x_i$, $i=0,\ldots,n$, car $f(x_i)=p_n(x_i)$ et $\pi_{n+1}(x_i) = 0$. Donc:
    \[
    \varphi(x_i) = f(x_i) - p_n(x_i) - \frac{f(x)-p_n(x)}{\pi_{n+1}(x)} \pi_{n+1} (x_i) = 0.
    \]
    De plus, elle s'annule aussi pour $t=x$ :
    \begin{eqnarray*}
    && \varphi(x) \\&=& f(x) - p_n(x) - \frac{f(x)-p_n(x)}{\pi_{n+1}(x)} \pi_{n+1} (x)\\
    &=& f(x) - p_n(x) - ( f(x) - p_n(x) ) = 0.
    \end{eqnarray*}
    En conséquence, $\varphi$ s'annule en $n+2$ points sur l'intervalle $[a,b]$: par application du Théorème de Rolle, sa dérivée $\varphi'$ s'annule en $n+1$ points (on vérifie, bien sûr, que $\varphi$ est bien continue sur $[a,b]$, dérivable sur $]a,b[$ de dérivée continue).
    
    \item Si $n=0$, on a $p_0 = f(x_0)$, et il suffit d'appliquer le Théorème des Accroissements finis entre $x$ et $x_0$: il existe $\xi \in ]a,b[$ tel que 
    \[
    f(x) = f(x_0) + f'(\xi) ( x - x_0 ),
    \]
    et on conclut en observant que $\pi_1(x) = (x - x_0)$:
    \[
    f(x) = p_0(x) + f'(\xi) \pi_1(x),
    \]
    
    \item Dans le cas $n\geq1$ :
    \begin{enumerate}
        \item En répétant $n$ fois l'argument de la question 2, on en déduit que la dérivée $n+1$-ème de $\varphi$ s'annule une fois sur $]a;b[$, ce qui permet de conclure: il existe $\xi \in ]a,b[$ tel que
        \[
        \varphi^{(n+1)}(\xi) = 0.
        \]
        
        \item Explicitons la dérivée $n+1$-ème de $\varphi$:
        \[
        \varphi^{(n+1)} (t) = f^{(n+1)}(t) - \frac{f(x)-p_n(x)}{\pi_{n+1}(x)} \pi_{n+1}^{(n+1)} (t),
        \]
        car $p_n$ est un polynôme de degré maximal $n$, dont la dérivée $n+1$-ème est nulle. On a de plus :
        \[
        \pi_{n+1}^{(n+1)} (t) = (n+1)!
        \]
        car $\pi_{n+1}$ est un polynôme de degré $n+1$ et le coefficient devant le terme de degré maximal ($x^{n+1}$) est $1$.
        En combinant ces deux résultats, et en prenant $t=\xi$, on aboutit alors au résultat souhaité:
        \[
        0 = 
        \varphi^{(n+1)} (\xi) = f^{(n+1)}(\xi) - \frac{f(x)-p_n(x)}{\pi_{n+1}(x)} (n+1) !.
        \]
        
    \end{enumerate}
    
    \item L'égalité souhaitée est une conséquence directe des questions 3 (pour $n=0$) et 4 (pour $n\geq1$).
    En effet, pour $n=0$, on a 
    \[
    f(x) = p_0(x) + f'(\xi) \pi_1(x),
    \]
    Donc
    \[
    f(x) - p_0(x) = f'(\xi) \pi_1(x),
    \]
    ce qui est bien l'égalité que l'on veut pour $n=0$.
    Pour $n \geq 1$, on a
    \[
        0 = 
         f^{(n+1)}(\xi) - \frac{f(x)-p_n(x)}{\pi_{n+1}(x)} (n+1) !
    \]
    et donc
    \[
         f^{(n+1)}(\xi) = \frac{f(x)-p_n(x)}{\pi_{n+1}(x)} (n+1) !
    \]
    et puis
    \[
         \frac{\pi_{n+1}(x)}{(n+1) !}
         f^{(n+1)}(\xi) = f(x)-p_n(x).
    \]
    Ce qui est bien l'égalité souhaitée :
    \[
         f(x)-p_n(x) = \frac{\pi_{n+1}(x)}{(n+1) !}
         f^{(n+1)}(\xi).
    \]
    
    \item Il suffit de majorer la dérivée $n+1$-ème de $f$ (elle est bornée car cette dérivée est supposée continue). 
    \[
         |f(x)-p_n(x)| \leq \frac{|\pi_{n+1}(x)|}{(n+1) !}
         \max_{a< t <b} |f^{(n+1)} (t)|.
    \]
        On a alors l'inégalité souhaitée avec
    \[
    M_{n+1} = \max_{a< t <b} |f^{(n+1)} (t)|.
    \]
\end{enumerate}

{\bf Remarque :} en toute généralité, on ne peut pas montrer que le terme de droite dans l'expression ci-dessus tend vers $0$ quand $n$ devient grand. Il y a même des fonctions pour lesquelles ce n'est pas le cas (phénomène de Runge), par exemple:
\[
f(x) = \frac{1}{1+x^2}.
\]

\newpage

\section{TD 6}

{\bf Exercice 6.1}

\begin{enumerate}
    \item Le calcul est facile car on connaît la primitive de l'exponentielle (qui est elle-même !):
    \[
    \int_0^1 e^x \mathrm{d}x = [ e^x ]_0^1 = e^1 - e^0 
    = e - 1 \simeq 1.71828182845905.
    \]
    De même, on connait la primitive du cosinus, ce qui permet d'avoir :
    \[
    \int_0^\pi \cos(x) \mathrm{d}x = \left [ \sin(x) \right ]_0^\pi = \sin(\pi) - \sin(0) = 0 - 0 = 0. 
    \]
    
    \item Il n'y a pas de technique connue pour déterminer ces intégrales, en particulier pas de primitives.
    
    \item Il n'y a pas de formule analytique connue pour la primitive ici (!), ni de technique connue pour déterminer la valeur de l'intégrale.
    
\end{enumerate}

{\bf Exercice 6.2}

\begin{enumerate}
    \item Il suffit d'injecter d'abord formule d'interpolation pour $p_1$ :
    \[
    \int_a^b p_1(x) \mathrm{d}x =
    \int_a^b \left ( L_0(x) f(a) + L_1(x) f(b) \right ) \mathrm{d}x. % =
    %f(a) \int_a^b L_0(x) \mathrm{d}x + f(b) \int_a^b L_1(x) \mathrm{d}x.
    \]
    Puis d'utiliser la linéarité de l'intégrale, et on obtient :
    \[
    \int_a^b p_1(x) \mathrm{d}x =
    %\int_a^b \left ( L_0(x) f(a) + L_1(x) f(b) \right ) \mathrm{d}x =
    f(a) \int_a^b L_0(x) \mathrm{d}x + f(b) \int_a^b L_1(x) \mathrm{d}x.
    \]
    Reste alors à calculer :
    \[
    \int_a^b L_0(x) \mathrm{d}x = \frac12(b-a).
    \]
    Cela correspond tout simplement à l'aire d'un triangle rectangle : la fonction $L_0$ est une fonction affine, et l'aire sous la courbe correspondante est l'aire du triangle rectangle de base $(b-a)$ et de hauteur 1. Le même résultat est obtenu pour l'autre intégrale, ce qui permet d'obtenir la formule souhaitée :
    \[
    \int_a^b p_1(x) \mathrm{d}x =
    %\int_a^b \left ( L_0(x) f(a) + L_1(x) f(b) \right ) \mathrm{d}x =
    \frac12 p_1(a) + \frac12 p_1(b).
    \]
    
    \item On utilise la formule (5) :
    
    Pour la première intégrale, cela donne : 
    \[
    \int_0^1 e^x \mathrm{d}x \simeq 
    \frac12 ( e^0 + e^1 ) = \frac12 ( 1 + e ).
    \]
    Et on obtient la valeur approchée :
    \[
    \int_0^1 e^x \mathrm{d}x \simeq 1.85914091422952.
    \]
    Ce qui correspond à une erreur de : $1.40859085770477.10^{-1}$.
    
    Pour la deuxième intégrale on obtient, de la même façon :
    \[
    \int_0^\pi \cos(x) \mathrm{d}x 
    \simeq \frac12 (\cos(0) + \cos(\pi)) = 0. 
    \]
    Ce qui correspond à une erreur de  $0$ (!).

Pour la troisième intégrale, on a :
    \[
    \int_0^1 e^{x^2} \mathrm{d}x 
    \simeq \frac12 ( e^0 + e^{1^2} )
    \simeq 1.85914091422952.
    \]
Et pour la quatrième :
    \[
    \int_0^\pi \cos(x^2) \mathrm{d}x \simeq 
    \frac12 ( \cos(0) + \cos(\pi^2) )
    \simeq
    0.152861476018906. 
    \]

\end{enumerate}

{\bf Exercice 6.3}

\begin{enumerate}
    \item Calculons :
    \begin{eqnarray*}
    & & \int_a^b p_2(x) \mathrm{d}x\\
    & = &
    %\int_a^b \left ( L_0(x) f(a) + L_1(x) f(b) \right ) \mathrm{d}x =
    f(a) \int_a^b L_0(x) \mathrm{d}x 
    + f(m) \int_a^b L_1(x) \mathrm{d}x
    + f(b) \int_a^b L_2(x) \mathrm{d}x \\
    &=& \frac{(b-a)}{6} f(a)
    + \frac{4(b-a)}{6} f(m)
    + \frac{(b-a)}{6} f(b).
    \end{eqnarray*}
    ce qui est le résultat attendu.
    
    Pour le calcul des intégrales, on peut faire via changement de variable. Par exemple pour la première, on a :
    \[
    \int_a^b L_0(x) \mathrm{d}x
    = \frac{(b-a)}{2}
    \int_{-1}^1 \widetilde L_0 ( t) \mathrm{d}t 
    \]
    en utilisant le changement de variable
    \[
    x = m + \frac{(b-a)}{2} t,
    \]
    et avec la notation $\widetilde L_0 ( t) = L_0 ( m + \frac{(b-a)}{2} t ) $.
    Et puis
    \[
    \int_{-1}^1 \widetilde L_0 ( t) \mathrm{d}t 
    =
    \frac12 \int_{-1}^1 (t^2 - t) \mathrm{d}t = \frac13.
    \]
    Ce qui donne bien la valeur $\frac{(b-a)}{6}$ pour le premier terme. On peut faire de même pour les deux autres intégrales.
    
    \item Avec la formule (6) on obtient :

    Pour la première intégrale :
    \[
    \int_0^1 e^x \mathrm{d}x 
    \simeq
    \frac16 e^0 + \frac46 e^{\frac12} + \frac16 e^1
    \simeq 1.71886115187659.
    \]
    Ce qui correspond à une erreur de : $5.79323417547961.10^{-4}$.
    
    Pour la deuxième on procède à l'identique et on obtient :
    \[
    \int_0^\pi \cos(x) \mathrm{d}x \simeq 1.11022302462516.10^{-16}. 
    \]
    Ce qui correspond à une erreur de : $1.11022302462516.10^{-16}$.


    Et puis pour la troisième : 
    \[
    \int_0^1 e^{x^2} \mathrm{d}x \simeq 1.47573058253500.
    \]
    Et enfin pour la quatrième : 
    \[
    \int_0^\pi \cos(x^2) \mathrm{d}x \simeq -1.585212535427892. 
    \]

\end{enumerate}

{\bf Exercice 6.4}

\begin{enumerate}
    \item Il suffit d'écrire $p_n$ dans la base de Lagrange :
    \[
    p_n = \sum_{k=0}^n f(x_k) L_k,
    \]
    puis de calculer :
    \[
        \int_a^b p_n(x) \mathrm{d}x = 
        \int_a^b \left ( \sum_{k=0}^n f(x_k) L_k(x) \right ) \mathrm{d}x =
        \sum_{k=0}^n \left ( \int_a^b L_k(x) \mathrm{d}x \right ) f(x_k) 
    \]
    et on obtient bien la formule
    \[
    \int_a^b p_n(x) \mathrm{d}x =
    \sum_{k=0}^n w_k f(x_k)
    \]
    avec
    \[
    w_k = \int_a^b L_k(x) \mathrm{d}x.
    \]
    \item La majoration est obtenue directement à partir de l'estimation sur l'inter\-polée de Lagrange de l'exercice 5.3.
    En effet, en utilisant la question 1, on peut reformuler l'erreur de quadrature comme suit :
    \[
    E_n (f) = \int_a^b f(x) \mathrm{d}x - \int_a^b p_n(x) \mathrm{d}x
    = \int_a^b (f(x) - p_n(x)) \mathrm{d}x.
    \]
    Puis on majore l'expression en valeur absolue et on utilise l'estimation de l'exercice 5.3 :
    \[
    | E_n (f) | % = \int_a^b f(x) \mathrm{d}x - \int_a^b p_n(x) \mathrm{d}x
    \leq \int_a^b |f(x) - p_n(x)| \mathrm{d}x
    \leq \int_a^b \left ( \frac{M_{n+1}}{(n+1)!} | \pi_{n+1}(x) | \right ) \mathrm{d}x.
    \]
    On obtient bien, finalement :
    \[
    | E_n (f) | % = \int_a^b f(x) \mathrm{d}x - \int_a^b p_n(x) \mathrm{d}x
    %\leq \int_a^b |f(x) - p_n(x)| \mathrm{d}x
    \leq \frac{M_{n+1}}{(n+1)!} \int_a^b  | \pi_{n+1}(x) | \mathrm{d}x.
    \]
    \item Il suffit d'estimer
    \[
    \frac{M_2}{2!} \int_a^b | \pi_2(x) | \mathrm{d}x
    \]
    Et donc, en particulier :
    \[
    \int_a^b | \pi_2(x) | \mathrm{d}x
    = \int_a^b | ( x - a ) ( x - b) | \mathrm{d}x
    = - \int_a^b ( x - a ) ( x - b) \mathrm{d}x
    \]
    où on a utilisé le fait que l'expression à intégrer est tout le temps négative sur l'intervalle $[a,b]$.
    On refait le coup du changement de variable, et on introduit 
    \[
    x = m + \ell t,
    \]
    avec $m=\frac{a+b}{2}$ et $\ell=\frac{b-a}{2}$.
    Ceci permet d'écrire :
    \[
    \int_a^b | \pi_2(x) | \mathrm{d}x
    = - \ell \int_{-1}^1 ( m + \ell t - a ) ( m + \ell t - b) \mathrm{d}t
     = - \ell^3 \int_{-1}^1 ( 1 + t ) ( -1 +  t ) \mathrm{d}t
    \]
    où on a utilisé $m-a=\ell$ et $b-m=\ell$. Reste à terminer le calcul :
    \[
    \int_{-1}^1 ( 1 + t ) ( -1 +  t ) \mathrm{d}t
    =
     \int_{-1}^1 ( -1 + t^2 ) \mathrm{d}t
    = 
     \left [ -t + \frac{t^3}{3}\right ]_{-1}^1
    \]
    et donc
    \[
    \int_{-1}^1 ( 1 + t ) ( -1 +  t ) \mathrm{d}t
    =
     ( -1 + \frac13 ) - ( 1 - \frac13 )
    = 
      -2 + \frac23 = - \frac43.
    \]
    Ainsi, on a déjà :
    \[
    \int_a^b | \pi_2(x) | \mathrm{d}x
     = - \ell^3 \left ( - \frac43 \right )
     = \frac{4 \ell^3}{3}
     = \frac{4 (b-a)^3}{3 \times 8}
     = \frac{(b-a)^3}{6}.
    \]
    Et on obtient finalement :
    \[
    \frac{M_2}{2!} \int_a^b | \pi_2(x) | \mathrm{d}x
    = \frac{M_2}{12} (b-a)^3.
    \]
    \item De même, un calcul un peu long, mais très divertissant, permet d'aboutir à l'égalité ci-dessous :
    \[
    \frac{M_3}{3!} \int_a^b | \pi_3(x) | \mathrm{d}x
    = \frac{M_3}{192} (b-a)^4.
    \]
    En effet, estimons 
    \[
    \int_a^b | \pi_3(x) | \mathrm{d}x
    = 
    \int_a^b | ( x - a ) ( x - m ) ( x - b) | \mathrm{d}x
    \]
    et après changement de variable
    \[
    \int_a^b | \pi_3(x) | \mathrm{d}x
    = 
    \ell 
    \int_{-1}^1 | ( m + \ell t - a ) ( m + \ell t - m ) ( m + \ell t - b ) | \mathrm{d}t.
    \]
    On utilise ensuite, de nouveau, les relations $d=m-a=b-m$, et on a donc :
    \[
    \int_a^b | \pi_3(x) | \mathrm{d}x
    = 
    \ell^4 
    \int_{-1}^1 | ( 1 +  t  ) t ( t - 1 ) | \mathrm{d}t.
    \]
    On termine le calcul en utilisant un argument de symétrie (parité), en levant la valeur absolue et par intégration directe :
    \[
    \int_{-1}^1 | ( 1 +  t  ) t ( t - 1 ) | \mathrm{d}t
    = 
    2 \int_{0}^1 ( 1 +  t  ) t ( 1 - t )  \mathrm{d}t = \frac12.
    \]
    Pour trouver la primitive, on a utilisé $( 1 +  t  ) t ( 1 - t ) = (1-t^2) t = t - t^3$.
    Finalement :
     \[
    \int_a^b | \pi_3(x) | \mathrm{d}x
    = 
    \frac{\ell^4}2 = \frac{(b-a)^4}{2^5}
    = \frac{(b-a)^4}{32}.
    \]
    Et donc
    \[
    \frac{M_3}{3!} \int_a^b | \pi_3(x) | \mathrm{d}x = 
    \frac{M_3}{6 \times 32} (b-a)^4
    = \frac{M_3}{192} (b-a)^4.
    \]
\end{enumerate}

\newpage

\section{TD 7}

{\bf Exercice 7.1}

La formule (7) s'obtient directement à l'aide de la formule (5) du TD 6 appliquée à chaque sous-intervalle, et en regroupant les termes identiques. En détails, cela donne d'abord :
\[
\int_a^b f(x) \mathrm{d}x = \sum_{i=0}^{m-1} \int_{x_i}^{x_{i+1}} f(x) \mathrm{d}x
\]
puis on approche sur chaque sous-intervalle :
\[
\int_{x_i}^{x_{i+1}} f(x) \mathrm{d}x \simeq \frac{x_{i+1}-x_i}{2} (f(x_i) + f(x_{i+1}))
= \frac{h}{2} (f(x_i) + f(x_{i+1})).
\]
On obtient donc en sommant :
\[
\int_a^b f(x) \mathrm{d}x \simeq \frac{h}{2} \sum_{i=0}^{m-1} (f(x_i) + f(x_{i+1})).
\]
Si on détaille :
\[
\int_a^b f(x) \mathrm{d}x \simeq \frac{h}{2} 
( f(x_0) + f(x_1)) + (f(x_1) + f(x_2)) + \ldots + (f(x_{m-1}) + f(x_m)).
\]
On s'aperçoit alors que tous les termes pour les indices $i=1$ à $i=m-1$ se regroupent deux par deux. Ce qui donne bien :
\[
\int_a^b f(x) \mathrm{d}x \simeq h
( \frac12 f(x_0) + f(x_1) + \ldots + f(x_{m-1}) + \frac12 f(x_m)).
\]
\medskip

{\bf Exercice 7.2}

Pour la première intégrale :
    \[
    \int_0^1 e^x \mathrm{d}x.
    \]
On se souvient que la valeur exacte  est : $1.71828182845905$.

Voici ce que donne la méthode des trapèzes (colonne 1 : nombre de subdivisions ; colonne 2 : valeur obtenue ; colonne 3 : erreur d'intégration).

\begin{verbatim}
1      1.85914091422952   1.40859085770477e-01   
2      1.75393109246483   3.56492640057802e-02
4      1.72722190455752   8.94007609847147e-03
8      1.72051859216430   2.23676370525694e-03
1000   1.71828197164920   1.43190149959338e-07
\end{verbatim}

De même pour :
    \[
    \int_0^\pi \cos(x) \mathrm{d}x.
    \]
On obtient :
\begin{verbatim}
1       0.00000000000000e+00   0.00000000000000e+00   
2       1.11022302462516e-16   1.11022302462516e-16   
4       0.00000000000000e+00   0.00000000000000e+00
8       5.55111512312578e-17   5.55111512312578e-17    
1000   -6.48027638494586e-15   6.48027638494586e-15
\end{verbatim}

Puis pour :
    \[
    \int_0^1 e^{x^2} \mathrm{d}x.
    \]
On a :
\begin{verbatim}
1      1.85914091422952   
2      1.57158316545863   
4      1.49067886169886
8      1.46971227642967   
1000   1.46265219895408
\end{verbatim}

Et enfin pour :
    \[
    \int_0^\pi \cos(x^2) \mathrm{d}x. 
    \]
On a :
\begin{verbatim}
1      0.152861476018906  
2     -1.150694032566193   
4      0.649760900266590
8      0.602621301755942   
1000   0.565695737289673
\end{verbatim}

{\bf Exercice 7.3}

On obtient le résultat souhaité directement à l'aide de l'estimation obtenue à l'exercice 6.4, et par sommation. En effet, il suffit d'abord de reformuler $\mathcal{E}_1$, en remarquant :
\[
\mathcal{E}_1 = \int_a^b f(x) \mathrm{d}x - \sum_{i=0}^{m-1} \int_{x_i}^{x_{i+1}} p_1^i(x) \mathrm{d}x,
\]
où $p_1^i$ est le polynôme d'interpolation de Lagrange associé à l'intervalle $[x_i;x_{i+1}]$.
Et donc en utilisant la relation de Chasles pour scinder la première intégrale :
\[
\mathcal{E}_1 = \sum_{i=0}^{m-1} \left (\int_{x_i}^{x_{i+1}} f(x) \mathrm{d}x - \int_{x_i}^{x_{i+1}} p_1^i(x) \mathrm{d}x \right ).
\]
On retrouve le terme d'erreur de quadrature présenté au TD 6, qu'on peut alors majorer. On utilise d'abord une inégalité triangulaire :
\[
| \mathcal{E}_1 | \leq  \sum_{i=0}^{m-1} \left  | \int_{x_i}^{x_{i+1}} f(x) \mathrm{d}x - \int_{x_i}^{x_{i+1}} p_1^i(x) \mathrm{d}x \right |.
\]
Et puis on applique la majoration de l'exercice 6.4 sur chaque sous-intervalle :
\[
| \mathcal{E}_1 | \leq  \sum_{i=0}^{m-1} \left  ( \frac{M_2^i}{12} (x_{i+1}-x_i)^3 \right ).
\]
On utilise ensuite l'égalité $x_{i+1}-x_i=h$ et on majore :
\[
M_2^i = \max_{\zeta \in [x_i;x_{i+1}]} |f''(\zeta)| \leq
\max_{\zeta \in [a;b]} |f''(\zeta)|,
\]
cette dernière expression ne dépend plus alors de $i$, et on peut la sortir de la somme, et on obtient donc :
\[
| \mathcal{E}_1 | \leq  \left  ( \frac{1}{12} \max_{\zeta \in [a;b]} |f''(\zeta)|  \right ) \sum_{i=0}^{m-1} h^3.
\]
On peut finalement simplifier, en utilisant la relation $m h = b - a$ :
\[
| \mathcal{E}_1 | \leq  \left  ( \frac{1}{12} \max_{\zeta \in [a;b]} |f''(\zeta)|  \right ) m h^3
= \frac{(b-a)^3}{12 m^2} \max_{\zeta \in [a;b]} |f''(\zeta)|. 
\]
Ce qui est la formule que l'on souhaite obtenir.

\medskip

On remarque pour terminer que ceci permet d'établir la convergence de la méthode des trapèzes : quand $m$ tend vers l'infini (ou $h$ tend vers $0$), la valeur approchée de l'intégrale tend vers la valeur exacte. En plus on voit que la convergence est quadratique : l'erreur de quadrature diminue en $m^2$ (elle est divisée par 4 si on divise $m$ par 2).

\medskip

{\bf Exercice 7.4}

C'est une application directe de la formule (6) du TD 6 (on procède exactement comme à l'exercice 7.1).

\medskip

{\bf Exercice 7.5}

Pour la première intégrale,
    \[
    \int_0^1 e^x \mathrm{d}x.
    \]
On obtient avec Simpson :    
\begin{verbatim}
1      1.71886115187659   5.79323417547961e-04   
2      1.71831884192175   3.70134627019070e-05   
4      1.71828415469990   2.32624085172439e-06   
8      1.71828197405189   1.45592846445552e-07   
1000   1.71828182845901   3.19744231092045e-14
\end{verbatim}

Pour
    \[
    \int_0^\pi \cos(x) \mathrm{d}x.
    \]
On obtient ensuite :
\begin{verbatim}
1      1.11022302462516e-16   1.11022302462516e-16   
2      5.55111512312578e-17   5.55111512312578e-17   
4      8.32667268468867e-17   8.32667268468867e-17
8      1.52655665885959e-16   1.52655665885959e-16   
1000   6.81898114363033e-14   6.81898114363033e-14
\end{verbatim}

Pour :
    \[
    \int_0^1 e^{x^2} \mathrm{d}x.
    \]
On obtient également :
\begin{verbatim}
1      1.47573058253500   
2      1.46371076044560   
4      1.46272341467327
8      1.46265632138942   
1000   1.46265174590715
\end{verbatim}

Pour : 
    \[
    \int_0^\pi \cos(x^2) \mathrm{d}x. 
    \]
On obtient finalement :
\begin{verbatim}
1    -1.585212535427892   
2     1.249912544544184   
4     0.586908102252393
8     0.565164272233738   
1000  0.565693513604310
\end{verbatim}

{\bf Exercice 7.6}

On obtient le résultat souhaité directement à l'aide de l'estimation obtenue à l'exercice 6.4, et par sommation. On va détailler en reprenant exactement le même raisonnement que dans l'exercice 7.3.

\medskip

Il suffit d'abord de reformuler $\mathcal{E}_2$, en remarquant :
\[
\mathcal{E}_2 = \int_a^b f(x) \mathrm{d}x - \sum_{i=0}^{m-1} \int_{x_i}^{x_{i+1}} p_2^i(x) \mathrm{d}x,
\]
où $p_2^i$ est le polynôme d'interpolation de Lagrange, de degré 2 cette fois, associé à l'intervalle $[x_i;x_{i+1}]$.
Et donc en utilisant la relation de Chasles pour scinder la première intégrale :
\[
\mathcal{E}_2 = \sum_{i=0}^{m-1} \left (\int_{x_i}^{x_{i+1}} f(x) \mathrm{d}x - \int_{x_i}^{x_{i+1}} p_2^i(x) \mathrm{d}x \right ).
\]
On retrouve le terme d'erreur de quadrature présenté au TD 6, qu'on peut alors majorer. On utilise d'abord une inégalité triangulaire :
\[
| \mathcal{E}_2 | \leq  \sum_{i=0}^{m-1} \left  | \int_{x_i}^{x_{i+1}} f(x) \mathrm{d}x - \int_{x_i}^{x_{i+1}} p_2^i(x) \mathrm{d}x \right |.
\]
Et puis on applique la majoration de l'exercice 6.4 sur chaque sous-intervalle :
\[
| \mathcal{E}_2 | \leq  \sum_{i=0}^{m-1} \left  ( \frac{M_3^i}{192} (x_{i+1}-x_i)^4 \right ).
\]
On utilise ensuite l'égalité $x_{i+1}-x_i=h$ et on majore :
\[
M_3^i = \max_{\zeta \in [x_i;x_{i+1}]} |f^{(3)}(\zeta)| \leq
\max_{\zeta \in [a;b]} |f^{(3)}(\zeta)|,
\]
cette dernière expression ne dépend plus alors de $i$, et on peut la sortir de la somme, et on obtient donc :
\[
| \mathcal{E}_2 | \leq  \left  ( \frac{1}{192} \max_{\zeta \in [a;b]} |f^{(3)}(\zeta)|  \right ) \sum_{i=0}^{m-1} h^4.
\]
On peut finalement simplifier, en utilisant la relation $m h = b - a$ :
\[
| \mathcal{E}_2 | \leq  \left  ( \frac{1}{192} \max_{\zeta \in [a;b]} |f^{(3)}(\zeta)|  \right ) m h^4
= \frac{(b-a)^4}{192 m^3} \max_{\zeta \in [a;b]} |f^{(3)}(\zeta)|. 
\]
Ce qui est la formule que l'on souhaite obtenir.

\medskip

On remarque pour terminer que ceci permet à nouveau d'établir la convergence de la méthode de Simpson : quand $m$ tend vers l'infini (ou $h$ tend vers $0$), la valeur approchée de l'intégrale tend vers la valeur exacte. En plus on voit que la convergence est cubique : l'erreur de quadrature diminue en $m^3$ (elle est divisée par 8 si on divise $m$ par 2).

\medskip

{\bf Suggestion finale :} essayez de voir si on retrouve ce comportement sur les exercices 7.2 et 7.5. Que remarque-t'on pour la méthode des trapèzes ? et pour la méthode de Simpson ? Que peut-on en conclure ?

\end{document}
